\chapter{Cohomology and the Diagnostic Hierarchy}\label{ch:cohomology}

\begin{workplan}{\Recast}
\wpgoal{This is the chapter V2 changes most in \emph{role}. In V0 the
cohomology was a diagnostic hierarchy. In V2 the Hodge decomposition is the
routing law of the wake phase and the harmonic component is the growth
address --- so this chapter carries the architecture, not commentary on it.}

\wpsource{This chapter's existing text (\S``the hierarchy is vacuous as
currently fed'' is the key surviving content);
\texttt{MATH\_TOOLBOX.md}~III.1 (Proposition~8.2 and the reframe it forced),
III.2 (measuring a genuine $\eta$: four runs, one retraction), III.6 (the
$b_1$-swamping bug and the triangle rule);
\texttt{TATIANA\_V1\_IDENTITY\_MAP.md}~\S4.1, \S9.}

\wpsteps{
  \item \textbf{Open with Proposition~8.2, not with the hierarchy.}
    $[\cob x] = 0$ in $\Hone$ \emph{unconditionally} --- a coboundary is by
    definition not a gluing failure. This holds for every state, every sheaf,
    identity restrictions or not. It is trivial to prove and it is the deepest
    correct criticism the project received: the engine's only $1$-cochain is
    always exact, so growth has no address, structurally, no matter how the
    state is computed. \Proved. Everything else in the chapter is downstream
    of it.
  \item \textbf{State what it kills.} The original growth intuition (read the
    address off $\Hone$ of a state the system already computes) is dead
    permanently. So is the ``residuals of learned restriction maps'' repair,
    which is \Refuted\ for the identical reason --- those residuals are still
    coboundaries of the learned maps, still exact.
  \item \textbf{Give the fix as the reframe it is.} Growth needs a genuinely
    \emph{measured} $1$-cochain $\eta$ that is not the coboundary of any
    single global state, plus the Hodge split; the harmonic component is the
    address. State this as the pivot of the whole architecture.
  \item \textbf{Flag the V2-specific hole loudly.} In V0, $\eta$ came from
    organs emitting pairwise judgements that disagreed. \emph{V2 has no
    organs.} Where a genuinely non-exact $\eta$ comes from for a single kernel
    is \textbf{open} and is the most important unanswered question in the
    book. Do not paper over it; give it its own section and forward-reference
    Chapter~\ref{ch:open}.
  \item \textbf{Keep the elicitation lesson.} A symmetric
    $\{u,v,\text{agreement}\}$ schema \emph{cannot} be a $1$-cochain ---
    antisymmetrising a symmetric table gives $\eta \equiv 0$ identically. Any
    future source of $\eta$, organ-based or not, must be antisymmetric by
    construction. Two earlier designs failed informatively before one worked.
  \item \textbf{Keep the dimensionality trap.} A $1$-dimensional harmonic
    subspace can report \emph{that} an obstruction exists but never
    \emph{where} --- at $b_1 = 1$ every harmonic component is a multiple of one
    fixed vector, so the ``consistent growth address'' it reports is an
    artifact. Any configuration used for a real address needs $b_1 \ge 2$.
    This nearly became a reported finding and must stay visible.
  \item \textbf{Keep the triangle rule and the swamping bug.} Record a
    $2$-simplex exactly when its three concepts co-fire; the earlier claim
    that ``no triangles $\Rightarrow b_1 = 0$'' was backwards --- no triangles
    \emph{maximises} $b_1$, and inserting each assembly as a clique
    manufactures $(n-1)(n-2)/2$ artifact cycles that swamp the real address.
    Also record that $b_1 = E - V + b_0 - F$ is \emph{wrong} once triangles
    share edges; compute $b_1$ from ranks, always.
  \item \textbf{Correct the headline statistic.} On a sparse graph a high
    ``holes'' fraction is the \emph{null expectation}
    ($\mathbb{E} = b_1/\abs{E}$), not evidence --- measured twelve sigma below
    a null in one run. The usable output is per-cycle \emph{localisation} of
    harmonic mass, which is what Construction~6 actually consumes. Replace the
    fraction with localisation wherever the old text reports it.
}

\wpdone{A reader can prove Proposition~8.2 unaided, explain why it makes the
growth address the central problem, and state what would have to be true of a
source of $\eta$ for it to work --- including that V2 does not yet have one.}
\end{workplan}

\section{Two pathologies that must not be confused}

\begin{definition}[Sheaf cohomology in low degrees]\label{def:cohom}
With coboundaries $\cob^k:\Ck{k}\to\Ck{k+1}$,
\[
  \Hzero=\ker\cob^0\quad(\text{global sections}),
  \qquad
  \Hone=\ker\cob^1/\im\cob^0\quad(\text{gluing failure}).
\]
\end{definition}

\begin{intu}\Intu
A \emph{structural hole} --- $b_1\neq0$, a Betti number of the bare
shape --- means a \emph{binding is missing}: no theorem yet ties some
loop together. A \emph{gluing failure} --- $\Hone\neq0$ of the sheaf ---
means the data are \emph{contradictory}: the organs cannot be
reconciled. The first calls for growth; the second calls for repair.
Treating a contradiction as a hole makes you add structure when you
should resolve a conflict, and treating a hole as a contradiction makes
you argue when you should build.
\end{intu}

\begin{remark}[$b_1$ is earned here]\label{rem:b1earned}
$b_1$ is a defensible diagnostic \emph{only} because simplices were
defined to mean genuine $n$-ary binding
(Remark~\ref{rem:nary}). A filled $2$-simplex really asserts a joint
relation, so an unfilled $1$-cycle really is a missing joint relation.
Most applications of topological data analysis to knowledge graphs skip
this justification and compute $b_1$ of a complex whose $2$-cells mean
nothing in particular, at which point the number is an artefact of the
construction rather than a fact about the knowledge.
\end{remark}

\section{The hierarchy is vacuous as currently fed}

This is the most important negative result in the book, and it is worth
stating as a proposition rather than as a caveat.

\begin{proposition}[Exactness of the emitted cochain]\label{prop:vacuous}
\Proved\ Suppose the $1$-cochain submitted for cohomological analysis is
obtained as $y=\cob^0 x$ for some global state $x\in\Czero$. Then
$[y]=0$ in $\Hone$, for every $x$, every sheaf $\Fs$, and every choice
of restriction maps.
\end{proposition}
\begin{proof}
$\Hone=\ker\cob^1/\im\cob^0$ and $y\in\im\cob^0$ by construction, so
$y$ represents the zero class. (That $y\in\ker\cob^1$ is automatic from
$\cob^1\cob^0=0$, but is not even needed: membership in the image is
already the conclusion.)
\end{proof}

\begin{hon}\Hon
\Refuted\ The consequence is not subtle. The engine computes its
$1$-cochain by applying $\cob$ to the current state --- that is what
$\discord$ is built from. So the class it then measures in $\Hone$ is
identically zero, and the growth mechanism that was supposed to be
addressed by a nonzero class \emph{has no address}. Every reported
$\Hone$ diagnostic in that pipeline was measuring the definition of a
coboundary, not the state of the system.

The repair requires a source of $1$-cochains that is \emph{not} of the
form $\cob x$. The natural candidate is pairwise judgements formed
independently: an organ pair that renders a verdict on their shared
edge without either of them consulting a global state. That is new
engineering --- organs that emit pairwise judgements, storage for
$2$-simplices, and an actual $\cob^1$ --- and it does not exist yet.

One proposed repair in the sources is recorded by the audit as
\emph{half circular}, and the circularity must be located precisely
before it is adopted. See Open Problem~\ref{op:H1}.
\end{hon}

\begin{remark}[The negative result is publishable]
Proposition~\ref{prop:vacuous} is not merely an internal correction. It
identifies a failure mode that any sheaf-based diagnostic over a derived
cochain will share, and the derivation is three lines. It belongs in
the paper on cohomological obstruction as a stated negative result,
not in a footnote.
\end{remark}

\section{The Hodge decomposition}

\begin{theorem}[Hodge decomposition for cellular sheaves]\label{thm:hodge}
\Proved\ Let $\Lapk{1}=\cob^0(\cob^0)^\top+(\cob^1)^\top\cob^1$ be the
degree-one Laplacian. Then
\begin{equation}\label{eq:hodge}
  \Cone=\im\cob^0\ \oplus\ \ker\Lapk{1}\ \oplus\ \im(\cob^1)^\top,
\end{equation}
an orthogonal decomposition, and $\ker\Lapk{1}\cong\Hone$.
\end{theorem}
\begin{proof}
Orthogonality of $\im\cob^0$ and $\im(\cob^1)^\top$ follows from
$\cob^1\cob^0=0$: for $y=\cob^0a$ and $z=(\cob^1)^\top b$,
$\inner{y}{z}=\inner{\cob^0a}{(\cob^1)^\top b}
=\inner{\cob^1\cob^0a}{b}=0$.
For $w\in\ker\Lapk{1}$ we have
$0=\inner{\Lapk{1}w}{w}=\norm{(\cob^0)^\top w}^2+\norm{\cob^1w}^2$, so
$w\in\ker(\cob^0)^\top\cap\ker\cob^1$, which is the orthogonal
complement of $\im\cob^0\oplus\im(\cob^1)^\top$. Counting dimensions
gives \eqref{eq:hodge}. Finally each class in
$\Hone=\ker\cob^1/\im\cob^0$ has a unique representative orthogonal to
$\im\cob^0$, and that representative lies in $\ker\Lapk{1}$.
\end{proof}

\begin{intu}\Intu
The decomposition sorts every $1$-cochain into three kinds of thing:
the part that is the disagreement of some global assignment
($\im\cob^0$), the part that is a genuine irreducible contradiction
($\ker\Lapk{1}$, the harmonic part), and the part that is a local
circulation with no bearing on gluing. Proposition~\ref{prop:vacuous}
says that as currently fed, the engine only ever produces cochains of
the first kind.
\end{intu}

\begin{workedbox}\Worked
On the running example $\Kc_4$ with the trivial sheaf
($\stalk{\sigma}=\R$): $\dim\Czero=4$, $\dim\Cone=4$ and the complex has
one filled triangle so $\dim\Ctwo=1$. Since the $1$-skeleton is
connected, $b_0=1$ and $\rank\cob^0=4-1=3$. The only cycle in the
$1$-skeleton is the triangle $\{\mathsf S,\mathsf M,\mathsf R\}$, and it
is filled, so $b_1=0$. Consistently:
$\dim\ker\cob^1=\dim\Cone-\rank\cob^1=4-1=3=\rank\cob^0$, so
$\Hone=0$.

Now delete the $2$-simplex, leaving the same $1$-skeleton unfilled. Then
$\Ctwo=0$, $\ker\cob^1=\Cone$ has dimension $4$, and
$\Hone\cong\R^{4-3}=\R$, matching $b_1=1$: one unfilled loop, one
class. This is the sense in which coning (Chapter~\ref{ch:cone}) kills
a class --- it adds precisely the $2$-cells that make the cycle a
boundary.
\hfill$\square$
\end{workedbox}

\section{The uncertainty stack}\label{sec:uncert}

\begin{intu}\Intu
Coherence is not correctness. Five instruments, ordered by cost, answer
five different questions, and they are not interchangeable. The most
common error in systems of this kind is to let the cheapest instrument
speak for all five.
\end{intu}

\begin{center}
\renewcommand{\arraystretch}{1.3}
\begin{tabular}{@{}p{2.7cm}p{5.2cm}p{2.5cm}p{2.9cm}@{}}
\toprule
\textbf{Instrument} & \textbf{Question} & \textbf{Cost} & \textbf{When}\\
\midrule
$\rhoscore$ & Are my parts agreeing? & effectively free & every tick\\
$\hat\rhoscore$ vs $\rhoscore$ & Did I \emph{know} whether I would agree? & free ($k$-NN) & once history exists\\
$b_0,\ b_1$ & What is structurally missing? & medium & on conflict\\
perturbation & Am I stable to paraphrase? & $k$ model calls & high stakes only\\
\textsc{VerifyOp} & Am I actually \emph{right}? & medium & promotion gate\\
\bottomrule
\end{tabular}
\captionof{table}{The five instruments and their distinct questions.}
\end{center}

\begin{remark}[Perturbation is a variance, not a derivative]\label{rem:perturbation}
Stability is measured as paraphrase-perturbation variance: run $k$
semantically equivalent paraphrases and measure the spread of outcomes.
This replaces an earlier proposal to compute a Fr\'echet derivative
$\norm{\partial F/\partial x}$ of the whole pipeline. That derivative
does not exist: $F$ factors through discrete token sampling, a DAG
generator, and database lookups, none of which is differentiable.
Earlier drafts also posited an infinite-dimensional Banach state space
for the same purpose; it was removed for the same reason. The
replacement costs $k$ model calls and must be rationed.
\end{remark}

\begin{remark}[Expected precision and the feedback edge]\label{rem:rhohat}
$\hat\rhoscore$ is a similarity-weighted $k$-nearest-neighbour
prediction of $\rhoscore$ from a log of past (query embedding,
$\rhoscore$) pairs. The \emph{signed} gap $\hat\rhoscore-\rhoscore$
separates three states that a single number cannot: \emph{competence}
(predicted agreement, observed agreement), \emph{known ignorance}
(predicted disagreement, observed disagreement --- ``confidently, I do
not know this''), and \emph{hallucination risk} (predicted agreement,
observed disagreement). A \textsc{VerifyOp} failure updates the store,
so self-knowledge improves with use. This is the only instrument in the
stack whose accuracy is supposed to increase over time, and it is
therefore the one that carries the architectural bet of
Chapter~\ref{ch:limits}.
\end{remark}
